%! Linear Transformations — Unit 5, Lesson 5.3 — Homework
\documentclass[10pt]{article}
\usepackage{linalg-article}
\usepackage{linalg-boxes}
\newcommand{\TallMath}[1]{$\displaystyle #1\rule[-1.4em]{0pt}{3.2em}$}
\newcommand{\vv}{\mathbf{v}}
\newcommand{\ww}{\mathbf{w}}
\newcommand{\xx}{\mathbf{x}}
\newcommand{\ee}{\mathbf{e}}
\newcommand{\T}{\mathsf{T}}
\newcommand{\zero}{\mathbf{0}}

\begin{document}

\pageheader{Unit 5, Lesson 5.3}{Homework}
\namedateperiod

\vspace{0.12in}

\begin{notesbox}{Practice}
Read each matrix as a \textbf{motion of the plane}: the columns are where $\ee_1,\ee_2$ land, the area
factor is $|\det A|$, and the \emph{sign} of $\det A$ tells you whether orientation flipped.

\begin{enumerate}[leftmargin=1.9em, itemsep=11pt]
  \item \textbf{Image areas.} For each matrix, compute $\det A$ and the area of the image of the unit square.
    \begin{enumerate}[label=(\alph*), leftmargin=1.8em, itemsep=5pt]
      \item $\begin{bmatrix} 5 & 0 \\ 0 & 2 \end{bmatrix}$ \hfill
      \item $\begin{bmatrix} 1 & 3 \\ 0 & 1 \end{bmatrix}$ \; {\footnotesize(a shear)} \hfill
      \item $\begin{bmatrix} 2 & 1 \\ 1 & 2 \end{bmatrix}$
    \end{enumerate}
    \vspace{1.2cm}
  \item \textbf{Name the motion and the orientation.} For each matrix, name the motion and state whether
        orientation is preserved or flipped (use the sign of $\det$).
    \begin{enumerate}[label=(\alph*), leftmargin=1.8em, itemsep=5pt]
      \item $\begin{bmatrix} 1 & 0 \\ 0 & -1 \end{bmatrix}$ \hfill \blank{3.5cm}
      \item $\begin{bmatrix} 3 & 0 \\ 0 & 3 \end{bmatrix}$ \hfill \blank{3.5cm}
      \item $\begin{bmatrix} 0 & -1 \\ 1 & 0 \end{bmatrix}$ \hfill \blank{3.5cm}
    \end{enumerate}
    \vspace{0.3cm}
  \item \textbf{Composition.} Let $A=\begin{bmatrix} 2 & 0 \\ 0 & 3 \end{bmatrix}$ and
        $B=\begin{bmatrix} 1 & 1 \\ 0 & 1 \end{bmatrix}$. Compute $\det A$ and $\det B$. Form $AB$ (do $B$,
        then $A$), compute $\det(AB)$, and check that it equals $\det A\cdot\det B$. What area factor does
        the combined map apply? \writelines{3}
  \item \textbf{Explain.} The transformation of a matrix $A$ sends the whole plane onto a single line.
        What is $\det A$? Is $A$ invertible? Explain in one sentence. \writelines{2}
\end{enumerate}
\end{notesbox}

\begin{extensionbox}
\textbf{Rotations and volume.}
\begin{enumerate}[label=(\alph*), leftmargin=1.8em, itemsep=5pt]
  \item \textbf{Rotation.} The $90^\circ$ rotation $R=\begin{bsmallmatrix} 0 & -1 \\ 1 & 0 \end{bsmallmatrix}$
        turns the plane. Compute $\det R$, and explain why \emph{every} rotation must have determinant $1$
        (think about area and orientation).
  \item \textbf{Volume in 3-D.} $|\det A|$ is a \emph{volume} factor in space. For the diagonal
        $A=\begin{bsmallmatrix} 2 & 0 & 0 \\ 0 & 2 & 0 \\ 0 & 0 & 3 \end{bsmallmatrix}$, compute $\det A$
        and give the volume of the image of the unit cube. Then do the same for
        $\begin{bsmallmatrix} 1 & 0 & 0 \\ 0 & 1 & 0 \\ 0 & 0 & -1 \end{bsmallmatrix}$ and say what the sign means.
\end{enumerate}
\writelines{3}
\end{extensionbox}

\begin{spiralbox}
\textbf{Coming next --- Unit 6: Eigenvalues and Eigenvectors.} Under a transformation, most vectors get
\emph{rotated} into a new direction. But a few special directions are only \emph{stretched} --- the arrow
comes out longer or shorter but pointing the same way. Those directions (\emph{eigenvectors}) and their
stretch factors (\emph{eigenvalues}) are the key to understanding a transformation you apply over and over.
\end{spiralbox}

\end{document}
